\section{Why the tree amplification gives numerical growth}\label{sec:growth}
The tree formulas are elaborate, so their intended numerical effect needs
proof rather than an appeal to their resemblance to ordinal notation.
A useful abstraction separates this question from the details of reflection.

Let \(v:\W\to\N\) satisfy just
\begin{equation}\label{eq:v-hypotheses}
 v(a)<v(\Suc(a)),\qquad v(s(0))<v(\Lim(s)).
\end{equation}
For each fixed \(n\), the interpretation \(v(a)=H_a(n)\) satisfies
these conditions by Lemma~\ref{lem:local}. No ordering of arbitrary
trees, monotonicity of the oracle, or monotonicity of a limit's branches
is assumed.

\begin{theorem}[Checked amplification properties]\label{thm:amplification}
For every \(v\) satisfying \eqref{eq:v-hypotheses}, every function
\(f:\W\to\W\), and every \(a\in\W\),
\[
 v(a)<v(T_f(a)),\qquad v(a)<v(P(f)(a)),\qquad v(a)<v(G(a)).
\]
Moreover, \(v(\gamma)<v(\Omega_*)\), where \(\gamma\) is the tree called
\code{gammaTree}.
\end{theorem}
\begin{proof}
The zeroth branch in \(T_f(a)\), after the descent operation is unfolded,
is
\[
                  \Suc\bigl(\Lim(k\mapsto f^k(a))\bigr).
\]
The zeroth child of its limit is \(f^0(a)=a\). Two applications each
of successor and limit growth establish the first claim. This argument
does not require \(f\) itself to increase \(v\).

If \(v(x)\le v(h(x))\) for every \(x\), structural induction on \(t\)
shows
\[
                         v(a)\le v(I_h(t,a)).
\]
The zero case is equality. A successor applies extensivity of \(h\)
after the induction hypothesis. At a limit, use its zeroth branch
and the second hypothesis in \eqref{eq:v-hypotheses}.

Apply this observation to \(h=T_f\), which is extensive by the first
claim. The outer successor in
\(P(f)(a)=\Suc(I_{T_f}(T_f(a),a))\) makes the inequality strict.
For \(G\), its zeroth limit branch is \(P(f_0)(a)\); apply the second
claim, followed by limit and successor growth.

Finally let \(s=G(\gamma)\). The third claim gives \(v(\gamma)<v(s)\)
and makes \(G\) extensive. Consequently
\[
             v(\gamma)<v(s)\le v(I_G(s,s))<v(\Omega_*).
\]
\end{proof}
\checked{higherTower\_increases;\quad stageIter\_extensive}
\checked{promoteTreeFunction\_increases;\quad promotedTreeFunction\_increases}
\checked{finalStage\_exceeds\_gammaTree}

Specializing this theorem to \(v(a)=H_a(n)\) gives, for every input \(n\),
\[
 H_a(n)<H_{T_f(a)}(n),\quad H_a(n)<H_{P(f)(a)}(n),\quad
 H_a(n)<H_{G(a)}(n),\quad H_\gamma(n)<H_{\Omega_*}(n).
\]
These specializations are separate checked theorems in
\code{NumericTreeGrowth.lean}. They establish actual numerical growth.
They do not identify a standard fast-growing-hierarchy index, prove
that \(P(f)(a)\) dominates \(f(a)\), or compare each adjacent finite
functional level with the next.

\section{The comparison and a fivefold exponential margin}\label{sec:comparison}
\subsection{The short strict-comparison proof}
The final tree has an outer successor. Thus \(N=\RF(H_b)(43)\)
for its explicitly defined predecessor \(b\), and
Lemma~\ref{lem:baseline}, with \(r=42\), immediately yields \(C_5<N\).
The complete comparison theorem is:
\begin{LeanCode}
theorem contender_5_lt_contender_6 :
    contender_5 < contender_6 := by
  unfold contender_6 finalStage
  exact baseline_lt_reflect _ 42
\end{LeanCode}
This proof is short because a uniform comparison with every oracle has
already been proved. Checking the supporting definitions and lemmas is
included in the reported budgets; the timing does not merely measure these
three final lines.

\subsection{A finite certificate inside the enormous tree}
For a baseline depth \(r\), put
\[
 B_r(x)=2^{x+r+2},\qquad B_r^{\,0}(x)=x,\qquad
 B_r^{\,k+1}(x)=B_r(B_r^{\,k}(x)).
\]
The relation \(\mathsf{MarginPath}(k,a)\) records a finite route through
\(a\) containing \(k\) successor edges, with only zeroth-child moves
at limit nodes. It has the rules
\[
 \mathsf{MarginPath}(0,a),
 \qquad
 \frac{\mathsf{MarginPath}(k,a)}
      {\mathsf{MarginPath}(k+1,\Suc(a))},
 \qquad
 \frac{\mathsf{MarginPath}(k,s(0))}
      {\mathsf{MarginPath}(k,\Lim(s))}.
\]
The endpoint of the route may be any tree. We need neither unfold the
whole tree nor traverse every branch.

\begin{lemma}[Path bound]\label{lem:path}
If \(k\ne0\) and \(\mathsf{MarginPath}(k,a)\), then
\[
                         B_r^{\,k}(C(r))\le H_a(r+1).
\]
\end{lemma}
\begin{proof}
Induct on the path certificate. A path of length zero is excluded by
the hypothesis. At the innermost successor, use the uniform bound
\[
                    B_r(C(r))\le\RF(H_b)(r+1),
\]
which holds for every predecessor \(b\), regardless of its value.
At each remaining successor, combine the induction hypothesis,
monotonicity of the ordinary function \(B_r\), and
\[
                    B_r(H_b(r+1))\le H_{\Suc(b)}(r+1).
\]
At each zeroth limit move, the corresponding child's value is bounded
by the finite maximum defining its parent.
\end{proof}
\checked{margin\_path\_bound}

The following certificate checks by unfolding only the required outer
constructors:
\begin{LeanCode}
theorem finalStage_five_successors :
    MarginPath 5 finalStage := by
  apply MarginPath.next
  apply MarginPath.next
  apply MarginPath.limit0
  apply MarginPath.next
  apply MarginPath.next
  apply MarginPath.limit0
  apply MarginPath.next
  exact MarginPath.base _
\end{LeanCode}
The moves have the following provenance:
\begin{center}
\begin{tabular}{cl}
\toprule
Successor counted & Constructor exposed along the chosen route\\
\midrule
1 & The outer successor of \(\Omega_*\)\\
2 & The outer successor of \(G\), exposed by the tree iterator\\
3 & The outer successor of \(P(f_0)\), selected at branch zero of \(G\)\\
4 & The outer successor of \(T_{f_0}\), exposed by its iterator\\
5 & The successor in the level-zero amplifier, after descent\\
\bottomrule
\end{tabular}
\end{center}
Limit moves between these successors use only branch zero, which is
always among the sampled branches. This route certifies only a small
part of the definition's growth.

\begin{theorem}[Certified numerical margin]\label{thm:margin}
Define \(x_0=C_5\) and \(x_{j+1}=2^{x_j+44}\). Then
\[
                                \boxed{x_5\le N.}
\]
\end{theorem}
\begin{proof}
Apply Lemma~\ref{lem:path} to \(r=42\), \(k=5\), and the displayed
certificate.
\end{proof}
\begin{LeanCode}
theorem contender_6_fivefold_margin :
    marginTower 42 5 (baseline 42) <= contender_6 :=
  margin_path_bound 42 finalStage_five_successors
    (Nat.noConfusion)
\end{LeanCode}
Here \code{baseline 42} is exactly the defining body of
\code{contender\_5}. The formulation avoids needlessly asking the
elaborator to normalize the huge closed number.

The fivefold bound is a rigorous lower bound, not an estimate of \(N\),
an upper bound, or a claim that five exponentials capture its scale.
It is much stronger than the strict comparison required by the contest.

\section{Constructivity and the scope of the claims}\label{sec:scope}
\subsection{Why this is a constructive natural number}
Every operation in the submitted closed definition is built from
inductive constructors, total functions, structural recursors,
finite lists, finite maxima, and natural-number arithmetic.
Although \(\Lim\) stores a function with infinitely many possible
arguments, it is represented by a finite piece of executable code.
At a numerical limit stage, only finitely many arguments are requested.

A subtle point is that reflection may call its oracle at an input
larger than the current \(n\). This does not threaten termination:
the recursively available object is the entire function \(H_a:\N\to\N\)
at a structurally smaller stage. The recursion is on the stage,
not on a claim that every oracle call decreases its natural input.
The Lean kernel accepts the resulting recursor-based definition.

The submitted definition and proofs depend on no termination oracle,
choice principle, or semantic assertion that arbitrary programs halt. In particular,
the construction does not use Busy Beaver values. Finite syntactic
enumeration of a language with a total interpreter is sufficient.
The distinction between finite descriptions, enormous values, and
uncomputable growth is central to large-number contests; Aaronson's
exposition gives broader context~\cite{aaronson}.

Computable in principle does not mean feasible to evaluate. Neither
\(C_5\) nor \(N\) was expanded to a numeral, and none of the verification
commands attempts that calculation. The task is to check a finite
definition and finite proofs about its value.

\subsection{Trees are not automatically canonical ordinals}
The constructors of \(\W\) resemble zero, successor, and countable
limits. The seed constructions deliberately resemble familiar
ordinal operations. We nevertheless use raw trees throughout,
with no quotient by ordinal equality, no certified normal-form
predicate, and no proof that a branch sequence is cofinal in a
particular ordinal.

This distinction matters numerically. For example,
\[
 H_{\Z}(n)=n+1,\qquad
 H_{\Lim(k\mapsto\Z)}(n)=n+2.
\]
A conventional ordinal supremum would identify the second tree's
putative ordinal value with zero, but our numerical semantics
distinguishes the two presentations. Theorems about canonical
ordinal notation cannot simply be transferred to these trees.

Accordingly, \code{gammaTree} is a descriptive source name; this
article does not prove that it denotes \(\Gamma_0\).
We claim no exact Bachmann--Howard or other named ordinal calibration
for \(\Omega_*\), and no associated proof-theoretic lower bound.
Formal relationships between different constructive notions of
ordinals require real additional hypotheses and proofs; see
Kraus, Nordvall Forsberg, and Xu~\cite{constructive-ordinals}.
Our comparison and amplification theorems avoid that unproved step.

\subsection{Exactly what the diagonal theorem says}
Theorem~\ref{thm:diagonal} fixes one expression and one total oracle.
For every sufficiently large input, reflection exceeds that
expression evaluated with that input and oracle. The finite depth
of the fixed expression supplies the explicit threshold.
The theorem neither quantifies over every Lean program nor asserts
that all oracle programs are increasing.

In particular, pointwise \(f\le g\) does not by itself justify
\(\RF(f)\le\RF(g)\): programs may inspect or transform oracle outputs
in ways that fail to preserve order. None of our proofs requires
this inference. We use the explicit guard \(f(n)\), quotation
membership, and branches actually sampled by a finite maximum.

Similarly, the tree proofs establish the specific inequalities in
Theorem~\ref{thm:amplification}. Their strength is that they hold for
every supplied \(f\) without auxiliary monotonicity assumptions.
They should not be read as a classification of the entire hierarchy
by growth rate.

\subsection{What is checked, and what remains a trust boundary}
The Lean kernel checks the internal theorems. The source-to-source
correspondence with the original Coq baseline is a detailed manual
audit supported by exact source preservation and small execution
checks. A fully formal translation between Coq's and Lean's
foundations is not part of this deliverable.

An empty axiom list means the checked declaration has no
axiom dependencies in Lean's environment. It still presupposes
the ordinary correctness of the kernel and the environment used
to execute it. The package records the exact compiler version,
source hashes, commands, and audit output so that the result can
be checked independently on a normal installation.

\section{Reproducible verification}\label{sec:verification}
\subsection{A pinned, small dependency base}
The contest modules use Lean \(4.19.0\) and import only \code{Init}
from its standard distribution~\cite{lean-release}. All other
imports refer to the included source modules. The external
dependency-inspection tool additionally imports Lean's official
metaprogramming API to read compiled declarations. That tool is
part of the verification harness, not part of the submitted
number's definition.

The verifier compiles all project modules afresh in a temporary
directory, without using prebuilt project object files.
It records process wall time for every module. Its definition
budget conservatively includes the baseline and shared interpreter,
the oracle language, all tree definitions, and the hierarchy.
Its proof budget includes the comparison's supporting proofs,
the auxiliary theorems, the small checks, the generated axiom
audit, and the compiled dependency inspection.

\begin{center}
\begin{tabular}{@{}lr@{}}
\toprule
Verification item & Recorded result\\
\midrule
Definition modules, fresh aggregate & 1.404626 seconds\\
Proofs, examples, and both audits & 2.741057 seconds\\
Inventoried declarations with no axioms & 144\\
Constants reachable from the submitted value & 444\\
Previous contender reachable & No\\
Baseline maximum reachable & No\\
Small executable check groups & 9 of 9 matched\\
\bottomrule
\end{tabular}
\end{center}
The run completed on 13 October 2026 using the official Lean 4.19.0
release (commit \texttt{6caaee842e94}), Linux x86\_64 with glibc 2.39,
and Python 3.12.14. The recorded compiler executable SHA-256 is
\begin{center}
\scriptsize\texttt{92c3d35b5bfaa5e0fea413a775d504cf46cd95e1345df61c2274f76779e7e023}.
\end{center}
The complete JSON result includes each module's wall time and output,
all source hashes, the 144 axiom-audit results, and the compiled
project-dependency closure.


These are observed measurements on the recorded environment,
not a promise that any machine, under any load, achieves the same
times. They give substantial room within the adopted \(15\)-second
definition and \(60\)-second proof limits.

\subsection{Axiom and dependency checks}
The harness first rejects forbidden source constructs such as
\code{sorry} and \code{unsafe}, ignoring comments and string
literals. This is only an early diagnostic. Its decisive logical
check generates \code{\#print axioms} commands for the public
definitions, inductive declarations, and theorems and verifies that
each reported dependency list is empty. Critical results are
explicitly required, so a missing comparison or margin theorem
cannot accidentally pass as a vacuous audit.

The definition of the new number must also be independent of
the previous contestant constant. The harness traverses the
actual compiled constants reachable from
\code{Biggest.contender\_6}, including types and declaration
bodies. It confirms that neither \code{Biggest.contender\_5}
nor \code{Biggest.baseline} is reachable. The recorded closure
also omits \code{Biggest.largestOfDepth} and \code{Biggest.maxBy}.
The shared raw syntax, interpreter, and uniform enumerators
are intentionally reachable: they implement the language being
extended, not a use of the old winning number.

The distinction is semantic. Merely importing a source module
that also declares \code{contender\_5} does not make that
constant a dependency of the new value. Conversely, hiding
the constant behind a helper definition would still be found
by the compiled closure traversal.

\subsection{Small executable checks}
The verified sample records include the type and term counts
and baseline values in Section~\ref{sec:baseline}, expression
counts \(0,4,28,914\) at depths zero through three, raw-interpreter
samples, quotation preservation, oracle/input behavior, and
small reflection values. Their purpose is to catch porting and
wiring errors that readable examples expose well. They do not
replace the general enumeration and comparison proofs.

No high-depth term list or large tree value is evaluated.
The checker verifies the finite syntax of the submission and
its proof terms; this is precisely the resource being measured.

\subsection{Commands to reproduce}
Install the standard Lean \(4.19.0\) distribution and use
Python \(3\) with its standard library. From the archive root,
run:
\begin{LeanCode}
python3 verification/verify.py \
  --lean /path/to/lean-4.19.0/bin/lean \
  --source-dir lean \
  --output verification/reproduced-results.json
\end{LeanCode}
The \code{lean-toolchain} file also pins the release for an
ordinary Elan installation. The supplied verifier builds
dependencies in their correct order and does not require Lake
or Mathlib.

To rebuild the article, use a TeX installation containing
XeLaTeX, Latin Modern and DejaVu fonts, and the packages named
in the preamble:
\begin{LeanCode}
cd article
latexmk -xelatex -interaction=nonstopmode \
  -halt-on-error reflection_tree_contender.tex
\end{LeanCode}
The checked source distribution contains the main article,
its three included \code{.tex} files, the compiled PDF,
and the independent verification evidence.

The execution container needed one narrowly scoped compatibility
helper because its PID namespace and mounted procfs did not agree:
Lean's lookup of its own executable used a pathname unavailable
in that namespace. The included C source redirects exactly that
own-executable lookup to \code{/proc/self/exe}. The official
compiler and shared libraries were unchanged. This affects
executable discovery, not logical constants or kernel checking.
A normal installation does not need it; full details and binary
hashes are recorded in the runtime notes.

\section{Guide to the artifacts}
\begin{longtable}{@{}p{0.32\textwidth}p{0.62\textwidth}@{}}
\toprule
File or directory & Role\\
\midrule
\endhead
\code{lean/Baseline.lean} & Exact raw interpreter, enumeration, maximum, and the reference constant.\\
\code{lean/Enumeration.lean} & Soundness and completeness of the original finite enumerators.\\
\code{lean/Oracle.lean} & Quotation, input, oracle syntax, and the reflection operator.\\
\code{lean/OracleFacts.lean} & Uniform oracle growth and baseline comparison.\\
\code{lean/Diagonal.lean} & Exact expression enumeration and eventual diagonal domination.\\
\code{lean/Trees.lean} & All tree operations and the explicit final stage.\\
\code{lean/Hierarchy.lean} & Numerical tree interpretation and \code{contender\_6}.\\
\code{lean/TreeGrowth.lean} & Generic structural amplification proofs.\\
\code{lean/TreeMargin.lean} & Finite path certificate and the fivefold exponential bound.\\
\code{lean/NumericTreeGrowth.lean} & Growth theorems instantiated for the actual hierarchy.\\
\code{lean/Comparison.lean} & The strict comparison with \code{contender\_5}.\\
\code{lean/Checks.lean} & Small executable regression examples.\\
\code{lean/Main.lean} & Convenient import point and principal theorem audits.\\
\code{verification/} & Fresh-build script, measured results, logs, and runtime notes.\\
\code{reference/} & Frozen upstream files, their license, source hashes, and correspondence notes.\\
\code{article/} & Full article source and compiled PDF.\\
\bottomrule
\end{longtable}

\section{Directions for a subsequent contender}
The present construction is frozen after proof and review. Several
substantive extensions could support a later contest entry.

One route is to develop an explicit ordinal-notation invariant,
prove the required branch properties, and connect the resulting
numerical hierarchy to a standard growth classification. That
would make named ordinal comparisons justified rather than
suggestive. A separate route is to index the functional-type
construction by trees, using dependent products at limit types;
this would require a compatible descent and numerical
majorization theory.

Another improvement is to extend the total object language while
preserving a proved interpreter and finite, complete enumeration.
The reflection proof is modular: a language embedding that respects
values and has a controlled depth cost immediately supplies a
new comparison. A formal Coq-to-Lean simulation would strengthen
the current source-correspondence boundary independently of
making the number larger.

For the existing entry, richer finite path certificates could
prove stronger lower bounds without changing its definition.
They would refine what is known about the same \(N\), rather
than estimate its full value. Efficient duplicate elimination
in enumeration could also be justified by proving preservation
of the semantic maximum, although it is unnecessary for
definition checking.

There is no canonical last enlargement. The contribution here is
a precise closed term with an explicit improvement mechanism,
a strict comparison with the repository's actual baseline, and
a reproducible record of the proofs and resource measurements.

\clearpage
\begin{thebibliography}{9}
\raggedright
\bibitem{contest-rules}
Cody Roux and contributors.
\emph{Name the biggest number: contest rules.}
Repository README, frozen at commit
\texttt{d56dce013e6264292533efdc5e1de71863c6fb86}.
\url{https://github.com/codyroux/name-the-biggest-number}.

\bibitem{contender-source}
Cody Roux and contributors; Contender 5 by Samuel Gr\"utter.
\emph{Contender.v.}
Same frozen commit; an exact copy is included in this archive.
\url{https://github.com/codyroux/name-the-biggest-number/blob/d56dce013e6264292533efdc5e1de71863c6fb86/Contender.v}.

\bibitem{lean-axioms}
Lean development team.
\emph{The Lean Language Reference: Axioms.}
\url{https://lean-lang.org/doc/reference/latest/Axioms/}.
Accessed 12 October 2026; verification itself is pinned to 4.19.0.

\bibitem{lean-inductive}
Lean development team.
\emph{The Lean Language Reference: Inductive Types.}
\url{https://lean-lang.org/doc/reference/latest/The-Type-System/Inductive-Types/}.
Accessed 12 October 2026.

\bibitem{lean-release}
Lean development team.
\emph{Lean 4.19.0 release.}
\url{https://github.com/leanprover/lean4/releases/tag/v4.19.0}.

\bibitem{constructive-ordinals}
Nicolai Kraus, Fredrik Nordvall Forsberg, and Chuangjie Xu.
\emph{Type-Theoretic Approaches to Ordinals.}
2022. \url{https://arxiv.org/abs/2208.03844}.

\bibitem{aaronson}
Scott Aaronson.
\emph{Who Can Name the Bigger Number?}
1999. \url{https://www.scottaaronson.com/writings/bignumbers.html}.
\end{thebibliography}

